\documentclass[instructions.tex]{subfiles}
\begin{document}
 
\chapter{Il faut souvent demander à Dieu qu'il donne de bons pasteurs à son Église.}
 
\primo Quatre fois l'année (aux Quatre-Temps), l'Église jeûne et prie solennellement pour demander à Dieu d'envoyer de bons pasteurs à son peuple. Les fidèles doivent s'unir à ses intentions, offrir à Dieu de ferventes prières à cette fin.
 
Un pasteur exemplaire, qui inspire la vertu, empêche le vice, est un présent du ciel. Heureux les peuples qui ont un tel pasteur! Combien d'âmes conduira-t-il à Dieu! On doit le chérir, l'estimer, le respecter: malheur à ceux qui le persécutent! Un pasteur, au contraire, qui est sans zèle, qui ne veille pas sur sa paroisse, qui laisse la jeunesse à sa liberté, qui n'a pas soin de tirer les âmes du désordre, est un châtiment des péchés des peuples. Dieu leur donne de tels pasteurs dans sa colère, parce qu'ils n'en méritent pas d'autres.
 
Mais une paroisse eût-elle le malheur d'avoir un pasteur sans vigilance et sans vertu, les paroissiens ne seraient point en droit de vivre dans le vice, et de laisser vivre la jeunesse dans le libertinage. Au défaut des instructions d'un pasteur, n'ont-ils pas les lumières de leur conscience, les exemples des gens de bien, les livres de piété, les inspirations de Dieu? Ils sont toujours sans excuse, s'ils vivent mal.
 
\secundo Les personnes qui en ont le pouvoir, devraient employer leur crédit pour procurer à l'Église de saints pasteurs, faire élever des jeunes gens de mérite et de vertu, pour remplir un jour les emplois de l'Église. Si l'on observait cette conduite, si la cupidité des parens ou l'ambition des grands ne se mêlait jamais dans le choix des ministres de Dieu, on verrait encore aujourd'hui, dans le clergé, des Jérôme, des Ambroise, des Chrysostome.
 
O vous qui, par recommandation, par intérêt, ou par d'autres motifs humains, faites engager dans le clergé, dans le cloître, dans les bénéfices, dans les emplois de l'Église, des gens sans talens et sans vertu, quel tort ne leur faites-vous pas! quel tort ne faites-vous pas à l'Église! Vous leur faites tort; vous êtes la cause de leur perte. Il vaudrait mieux pour eux, dit saint Bernard, vivre en artisans, fouir la terre ou mendier, que d'être élevés à un état dont ils sont indignes, où Dieu ne les appelle pas\footnote{Bonum erat eis magis fodere aut mendicare. Bern.}.
 
Vous faites tort à l'Église, et un tort irréparable: en pensant lui donner des pasteurs et des ministres, vous donnez des persécuteurs à Jésus-Christ, qui scandalisent, qui laissent périr les fidèles\footnote{Ministri Christi sunt, et serviunt Antichristo. Bern.}; si, en contribuant par sa faute à la perte d'une seule âme on mérite l'enfer, combien d'enfers méritent ceux qui plaçant de mauvais pasteurs, contribuent à la perte de tant d'âmes!
 
Quant aux personnes qui donnent volontairement occasion de péché à des gens consacrés à Dieu, et qui sont complices de leurs fautes, combien ne sont-elles pas criminelles! Oh! qu'elles sont malheureuses d'être les instrumens du démon pour perdre les ministres de Dieu! Quel compte rendront-elles, au jour du jugement, des sacriléges qu'elles leur font commettre, des scandales qu'elles occasionnent, et de la perte des âmes dont elles sont la cause par ces chutes affreuses!
 
\section{Résolutions}
 
\primo Je prierai souvent en général pour tous les pasteurs de l'Église, et en particulier pour le mien.
 
\secundo Je contribuerai, selon mon pouvoir, à l'éducation des jeunes gens qui montreront des dispositions à l'état ecclésiastique, et le désir d'y entrer.
 
\tertio Et jamais aucune vue terrestre ne m'engagera à influer sur la vocation de mes enfans. 

\prayer{Mon Dieu, les besoins de votre Église sont grands: donnez-nous des pasteurs qui nous éclairent par leurs lumières, nous édifient par leurs vertus, et nous conduisent au ciel par la ferveur et la prudence de leur zèle.}
 
\end{document}
